\documentclass[11pt]{article}

\usepackage[a4paper,margin=2.4cm]{geometry}
\usepackage{amsmath,amssymb}
\usepackage{booktabs}
\usepackage{siunitx}
\usepackage[hidelinks]{hyperref}
\usepackage{enumitem}

\newcommand{\tirr}{t_{\mathrm{irr}}}
\newcommand{\tdel}{t_{\mathrm{del}}}
\newcommand{\tmeas}{t_{\mathrm{meas}}}

\title{Decay kinetics:\\
from produced emitters to measured decays, and $\sigma(\text{range})$}
\author{Method note}
\date{\today}

\begin{document}
\maketitle

\begin{abstract}
This note (\texttt{decay\_sampling/}) converts the proton run's production
of $\beta^{+}$ emitters into the annihilation source measured by a detector. It
scales the production to a clinical dose, applies radioactive decay and the
acquisition timing to obtain the measured number of decays per isotope, and
delivers the source as a budget. The figure of merit $\sigma(\text{range})$ is
the spread of the recovered proton range over Poisson realisations of that
budget.
\end{abstract}

\section{Role}

The proton run wrote \texttt{emitters.csv}: one row per produced emitter with its
annihilation point and isotope, i.e.\ the production shape $P_j(\mathbf r)$ at
the statistics of the simulated protons. The detector measures a different
quantity. This step applies two transforms: an absolute scale to a clinical
dose (Sec.~\ref{sec:norm}), and the decay-plus-window selection that yields the
measured decays (Sec.~\ref{sec:decay}).

\section{Time-decay model}
\label{sec:decay}

Species $j$ has decay constant $\lambda_j=\ln2/T_{1/2,j}$ and decays $100\%$ by
$\beta^{+}$. For constant production over $\tirr$ at rate $K_j=P_j/\tirr$, the
activity is
\begin{align}
A_j(t) &= K_j\bigl(1-e^{-\lambda_j t}\bigr), & 0\le t\le \tirr,\\
A_j(t) &= K_j\bigl(1-e^{-\lambda_j \tirr}\bigr)\,e^{-\lambda_j (t-\tirr)},
       & t>\tirr.
\end{align}
The scanner acquires over $[t_0,t_1]$ with $t_0=\tirr+\tdel$ and
$t_1=t_0+\tmeas$. The collected decays $N_j=\int_{t_0}^{t_1}A_j(t)\,dt$ evaluate
to the three-factor expression (continuous / cyclotron limit; acquisition is
post-beam, so the pulse structure does not enter):
\begin{equation}
N_j = P_j\;
\frac{1-e^{-\lambda_j \tirr}}{\lambda_j \tirr}\;
e^{-\lambda_j \tdel}\;
\bigl(1-e^{-\lambda_j \tmeas}\bigr).
\label{eq:nmeas}
\end{equation}
The three factors are the survival to the end of beam, the decay during the
transport delay $\tdel$, and the fraction of survivors that decay during the
acquisition $\tmeas$.

Operating point (in-room baseline): $\tirr=\SI{60}{s}$, $\tdel=\SI{120}{s}$,
$\tmeas=\SI{1200}{s}$ (cyclotron); all are parameters. $\tdel$ is the key swept
variable (it sets the surviving $^{15}$O); $\tdel=\SI{300}{s}$,
$\tmeas=\SI{1800}{s}$ is a conservative offline point.

\section{Absolute normalisation}
\label{sec:norm}

The proton run scores the target-box dose $D_{\mathrm{sim}}$ for the $N_{\mathrm{sim}}$
simulated protons (\texttt{run\_meta.csv}). The production for a clinical dose
$D$ is
\begin{equation}
P_j(D) = \mathrm{count}_j\,\frac{D}{D_{\mathrm{sim}}}.
\label{eq:pj}
\end{equation}
The Parodi Table~2 cross-check
(\texttt{analysis\_transport/parodi\_cross\_check.py}) validates
$P_j(\SI{1}{Gy})=\mathrm{count}_j/D_{\mathrm{sim}}$. With
Eqs.~\eqref{eq:pj},\eqref{eq:nmeas} the measured totals are
$N\sim10^{7}$--$10^{8}$.

\section{Source sampling}
\label{sec:source}

The detector requires $N_j$ annihilation points per isotope. They are drawn from
\texttt{emitters.csv} by sampling with replacement, reproducing the shape
$P_j(\mathbf r)$ at the required count. The handoff (this repo) writes only the
deterministic budget: the per-isotope counts $N_j$ and the operating point. The
Poisson realisations (Sec.~\ref{sec:sigma}) and the sampling of annihilation
points are done by the downstream detector study. The same budget and
seed yield an identical source for every detector.

\section{Figure of merit}
\label{sec:sigma}

$\sigma(\text{range})$ is the precision of the recovered proton range --- the
error bar on a single measurement. The realised count of species $j$ is
$M_j\sim\mathrm{Poisson}(N_j)$ about the expectation $N_j$ of
Eq.~\eqref{eq:nmeas}. One realisation $z$ draws $M_j^{(z)}$ for every isotope,
forms the source, runs the detector and the reconstruction, and fits the distal
fall-off to a range estimate $R_z$. Over $Z\sim100$ realisations,
\begin{equation}
\sigma(\text{range}) = \operatorname{std}\{R_1,\dots,R_Z\}.
\end{equation}
The detector with the smallest $\sigma(\text{range})$ at a fixed budget is
preferred. The detector is simulated once to a master coincidence list; the $Z$
realisations are Poisson resamples of that list and only the reconstruction
repeats (downstream detector study).

\section{Implementation (\texttt{decay\_sampling/})}

\begin{itemize}[nosep]
\item \texttt{budget.py} (deterministic): inputs \texttt{emitters.csv},
\texttt{run\_meta.csv}; parameters $D,\tirr,\tdel,\tmeas$. $\lambda_j$ from
\texttt{common/isotopes.py}; $N_j$ from Eqs.~\eqref{eq:pj},\eqref{eq:nmeas}.
Output \texttt{sampling\_budget\_<scenario>.csv} --- per isotope $N_j$ and the
operating point. No random numbers.
\item \texttt{budget\_gen.py} (stochastic): reads the budget, draws
$M_j^{(z)}\sim\mathrm{Poisson}(N_j)$ for $Z$ realisations with a seed; output
\texttt{sampling\_realizations\_<scenario>.csv}. This step moves to the
downstream detector study.
\end{itemize}

\end{document}
